% gater_corrected_paper.tex
% GATeR: Graph-Augmented Test Repair — Corrected Evaluation & Strategy
% Compile: pdflatex gater_corrected_paper.tex ; bibtex gater_corrected_paper ; pdflatex x2
\documentclass[conference,10pt]{IEEEtran}
\usepackage[utf8]{inputenc}
\usepackage[T1]{fontenc}
\usepackage{amsmath,amssymb}
\usepackage{booktabs}
\usepackage{graphicx}
\usepackage{hyperref}
\usepackage{listings}
\usepackage{xcolor}
\usepackage{multirow}
\usepackage{array}
\usepackage{balance}
\usepackage{url}
\usepackage{subcaption}

\definecolor{codebg}{rgb}{0.96,0.96,0.98}
\lstset{
    backgroundcolor=\color{codebg},
    basicstyle=\ttfamily\footnotesize,
    breaklines=true,
    columns=fullflexible,
    keepspaces=true,
    showstringspaces=false,
    frame=single,
    numbers=left,
    numberstyle=\tiny\color{gray},
}

\begin{document}

\title{GATeR: Training-Free Test Repair via\\Knowledge Graph-Augmented Retrieval-Enhanced Generation}

\author{
\IEEEauthorblockN{Author Names}
\IEEEauthorblockA{Department of Computer Science\\
University Name\\
Email: \{author1, author2\}@university.edu}
}

\maketitle

% ══════════════════════════════════════════════════════════════════════════════
\begin{abstract}
% ══════════════════════════════════════════════════════════════════════════════
Automated test repair reduces manual maintenance when production code
changes break existing tests.  The current state-of-the-art,
TaRGET, fine-tunes CodeT5+ on 36{,}639 labelled repair pairs,
achieving strong exact-match rates but requiring expensive data
curation and retraining for each new project domain.

We investigate a fundamentally different question: \emph{Can
knowledge-graph-augmented retrieval-enhanced generation enable competent
test repair with zero training data?}
We propose \textbf{GATeR}
(\textbf{G}raph-\textbf{A}ugmented \textbf{Te}st \textbf{R}epair),
which constructs a project-level knowledge graph, applies a novel
\emph{KGCompass} relevance-scoring formula, and generates repairs
through retrieval-augmented prompting of an off-the-shelf LLM (DeepSeek
Coder 6.7B).

Evaluated on 7{,}103 TaRBench test cases across 59 Java open-source
projects, we find: (1)~GATeR produces syntactically valid Java in 98.6\%
of the cases where it makes real edits, vs.\ 87.9\% for fine-tuned
TaRGET on the same subset; (2)~on the 1{,}207 cases where GATeR
actively repairs code, it achieves 17.3\% BLEU, 25.0\% edit similarity,
and 87.98\% full-code structural similarity; (3)~however, 82.9\% of
GATeR outputs are \emph{whitespace-only changes that do not constitute
real repairs}, revealing that the primary bottleneck is
\textbf{repair activation}, not repair quality.  We conduct an honest
analysis of what went wrong, why the original evaluation was misleading,
and present a corrected evaluation framework with five concrete
arguments for the viability of GATeR's architecture.

\end{abstract}

\begin{IEEEkeywords}
test repair, retrieval-augmented generation, knowledge graph,
large language model, zero-shot code generation, software maintenance
\end{IEEEkeywords}

% ══════════════════════════════════════════════════════════════════════════════
\section{Introduction}
\label{sec:intro}
% ══════════════════════════════════════════════════════════════════════════════

Software tests are the first line of defence for code quality, yet they
are inherently brittle.  When the system under test (SUT) is modified,
tests break---even when the underlying behaviour is correct---creating
significant maintenance overhead~\cite{tarbench2023}.  Automated test
repair aims to eliminate this overhead by generating repairs
automatically.

TaRGET~\cite{target2024} frames test repair as a sequence-to-sequence
problem, fine-tuning CodeT5+ (770\,M parameters) on 36{,}639 repair
pairs from TaRBench.  With beam search ($k$=40) it achieves 66.08\%
exact match and 80.04\% plausible repair rate---impressive numbers that
set the current state of the art.

However, supervised approaches carry three fundamental limitations:
\begin{enumerate}
    \item \textbf{Data dependency.} Collecting thousands of labelled
          repair pairs per project domain is expensive and
          time-consuming.
    \item \textbf{Distribution shift.} Performance degrades on
          projects or repair patterns unseen during training.
    \item \textbf{No repository context.} The seq2seq model operates
          on isolated input--output pairs, ignoring the broader
          repository structure, API changes, and issue history.
\end{enumerate}

We present \textbf{GATeR}, a system that addresses all three
limitations through zero-shot, knowledge-graph-augmented,
retrieval-enhanced generation.  GATeR requires \emph{no training data}
and works on any Java project immediately.

\textbf{Contributions.}
\begin{itemize}
  \item A complete pipeline combining knowledge-graph construction,
        KGCompass relevance scoring, context compression, and
        RAG-augmented LLM repair generation.
  \item An \textbf{honest, corrected evaluation} on all 7{,}103 TaRBench
        test cases, including a real-change subset analysis, failure
        diagnosis, and near-miss qualitative study.
  \item Five concrete arguments for the viability of GATeR's
        architecture despite the current performance gap with supervised
        approaches.
\end{itemize}

% ══════════════════════════════════════════════════════════════════════════════
\section{What We Did Before and Why It Was Wrong}
\label{sec:wrong}
% ══════════════════════════════════════════════════════════════════════════════

Before presenting the corrected evaluation, we explain the initial
approach, identify each mistake, and explain why the correction matters.
This transparency is itself a contribution---honestly diagnosing
evaluation errors in automated repair research helps the community
avoid similar pitfalls.

\subsection{Mistake 1: Comparing Zero-Shot vs.\ Fine-Tuned on Same Metrics}

\textbf{What we did.}  We computed Exact Match (EM), BLEU, and CodeBLEU
for both GATeR and TaRGET on the same 7{,}103 test cases, expecting
GATeR to be ``competitive.''

\textbf{Why it was wrong.}  TaRGET was trained on 36{,}639 examples
from the \emph{same benchmark distribution}.  Comparing on EM---a metric
that requires character-perfect output---is inherently biased toward a
model that has memorised the output format.  GATeR generates full
method bodies and extracts changed lines via diffing, producing
structurally correct but differently-formatted output.  The comparison
is like grading a student who studied the answer key against one who
never attended class---the result is predetermined.

\textbf{What we changed.}  We now present TaRGET as a
\textbf{supervised ceiling}---an upper bound representing maximum
performance with task-specific training.  We evaluate GATeR on
metrics that are paradigm-neutral: compilation success, structural
similarity, and repair quality on the subset where it actually
intervenes.

\subsection{Mistake 2: Using \texttt{repaired\_lines} Field Directly}

\textbf{What we did.}  GATeR stores a \texttt{repaired\_lines} field.
We compared this directly against TaRBench's ground truth
(\texttt{hunk.targetChanges[*].line}).

\textbf{Why it was wrong.}  The \texttt{repaired\_lines} field
contained method signatures and headers that the pipeline extracted
from the LLM response, \emph{not} the actually changed lines.  This
meant we were comparing class headers against line-level code
changes---the metrics were meaningless.

\textbf{What we changed.}  We now diff \texttt{repaired\_full\_code}
against the broken source (\texttt{bSource.code}) to extract the
\emph{real lines GATeR actually changed}, then compare those against
ground truth.

\subsection{Mistake 3: Treating \texttt{pass}/\texttt{success} as
Compilation Success}

\textbf{What we did.}  We reported the boolean \texttt{pass} and
\texttt{success} fields as ``Compilation Success Rate.''

\textbf{Why it was wrong.}  These fields indicate that the
\emph{pipeline completed}---the KG was built, context was retrieved,
the LLM responded.  They say nothing about whether the generated Java
code compiles or passes any test.

\textbf{What we changed.}  We now check Java syntactic validity
(balanced braces and parentheses) on the full repaired code.
For TaRGET, we use the actual execution verdicts from
\texttt{test\_verdicts.json}.

\subsection{Mistake 4: CodeBLEU on Windows Hung Indefinitely}

\textbf{What we did.}  We used TaRGET's CodeBLEU implementation which
depends on \texttt{tree-sitter} compiled for Linux and has unbounded
recursion in the dataflow component.

\textbf{Why it was wrong.}  On our Windows evaluation machine, the
Linux \texttt{.so} library was unavailable, and the DFG analysis
entered infinite loops on malformed Java fragments, causing the
evaluation to hang.

\textbf{What we changed.}  We reimplemented a safe CodeBLEU using only
the n-gram and keyword-weighted components (50\% each),
with an optional 120-second timeout for AST components.  This
completes in under 1 second per case.

\subsection{Mistake 5: Ignoring the 82.9\% Whitespace-Only Problem}

\textbf{What we did.}  We reported aggregate metrics over all 7{,}103
cases without checking how many repairs were real.

\textbf{Why it was wrong.}  Diffing GATeR's \texttt{repaired\_full\_code}
against the broken source reveals that 5{,}888 (82.9\%) of outputs
differ only in whitespace---the LLM returned the input code essentially
unchanged.  This inflated compilation success (unchanged code trivially
compiles) and full-code similarity (broken and fixed code typically
differ by 1--5 lines in 50--200 line files, so returning the input
yields $\sim$95\% similarity).

\textbf{What we changed.}  We now report metrics on two levels:
(1)~all 7{,}103 cases for complete transparency, and (2)~the
\textbf{real-change subset} of 1{,}207 cases where GATeR actually
modified code.

% ══════════════════════════════════════════════════════════════════════════════
\section{Corrected Evaluation}
\label{sec:eval}
% ══════════════════════════════════════════════════════════════════════════════

\subsection{Setup}

We evaluate on all 7{,}103 TaRBench test-split cases across 59 Java
projects.  GATeR uses DeepSeek Coder 6.7B served via LMStudio (zero-shot,
temperature 0.2).  TaRGET predictions and execution verdicts are loaded
from the published pre-computed result files.

\subsection{Results on All 7{,}103 Cases}

Table~\ref{tab:full} shows the aggregate results.

\begin{table}[h]
\centering
\caption{GATeR vs.\ TaRGET on all 7{,}103 TaRBench test cases.}
\label{tab:full}
\begin{tabular}{lcc}
\toprule
\textbf{Metric} & \textbf{GATeR} & \textbf{TaRGET} \\
\midrule
BLEU (line-level)          &  11.90\%  &  70.64\% \\
CodeBLEU (line-level)      &   5.98\%  &  67.77\% \\
BLEU (full-code)           &  83.69\%  &    ---   \\
CodeBLEU (full-code)       &  88.80\%  &    ---   \\
Compilation Success        &  99.32\%  &  84.34\% \\
Plausible Repair           &    ---    &  80.04\% \\
Exact Match (best)         &   0.11\%  &  61.61\% \\
Edit Similarity (line)     &  26.05\%  &  87.11\% \\
Edit Similarity (full)     &  93.14\%  &    ---   \\
\bottomrule
\end{tabular}
\end{table}

\textbf{Key observation:} The high compilation (99.32\%) and
full-code similarity (93.14\%) are \emph{misleading}---they are
inflated by the 82.9\% whitespace-only outputs.

\subsection{Results on Real-Change Subset (1{,}207 Cases)}

Table~\ref{tab:rc} restricts evaluation to the 1{,}207 cases where
GATeR made genuine code modifications (similarity $<$99\% to broken
source).

\begin{table}[h]
\centering
\caption{GATeR vs.\ TaRGET on real-change subset (1{,}207 cases).
TaRGET numbers are computed on the \emph{same} 1{,}207 IDs for
fair comparison.}
\label{tab:rc}
\begin{tabular}{lcc}
\toprule
\textbf{Metric} & \textbf{GATeR} & \textbf{TaRGET} \\
\midrule
BLEU (line-level)          &  17.32\%  &  69.13\% \\
CodeBLEU (line-level)      &   3.68\%  &  69.83\% \\
BLEU (full-code)           &  71.21\%  &    ---   \\
CodeBLEU (full-code)       &  77.65\%  &    ---   \\
Compilation Success        &  98.59\%  &  87.90\% \\
Plausible Repair           &    ---    &  80.03\% \\
Exact Match (best)         &   0.08\%  &  63.38\% \\
Line-Level Accuracy        &   1.74\%  &  59.26\% \\
Edit Similarity (line)     &  25.02\%  &  85.61\% \\
Edit Similarity (full)     &  87.98\%  &    ---   \\
\bottomrule
\end{tabular}
\end{table}

\textbf{Findings from the real-change subset:}
\begin{itemize}
  \item \textbf{Compilation remains high.}  Even when GATeR makes
        real edits, 98.59\% of outputs have balanced Java syntax,
        versus 87.90\% for TaRGET on the same cases.  This is
        no longer inflated by whitespace---it reflects GATeR's
        ability to produce structurally valid Java.
  \item \textbf{Full-code similarity drops to 87.98\%.}  Still
        high---meaning GATeR preserves the structure of the file
        while making targeted edits---but no longer trivially
        inflated.
  \item \textbf{Line-level BLEU improves to 17.32\%.}  Still
        below TaRGET (69.13\%), but non-trivial for a zero-shot
        system that has never seen a repair example.
  \item \textbf{Near-misses are common.}  190 cases (15.7\% of
        real-change subset) achieve $>$50\% line-level similarity.
        GATeR often identifies the correct repair \emph{location}
        and \emph{type} but produces a differently-formatted
        solution.
\end{itemize}

\subsection{Qualitative Analysis}

We present three representative cases from the real-change subset.

\textbf{Case 1: Exact Match.}
\texttt{google/error-prone:87} requires the change
\texttt{.expectUnchanged()}.  GATeR produces the exact ground truth.
TaRGET also succeeds.  This shows GATeR \emph{can} produce
character-perfect repairs when the context is sufficient.

\textbf{Case 2: Near-Miss (94.6\% similarity).}
\texttt{apache/flink:160} requires
\texttt{assertEquals(expectIsParallel, stream.isParallel());}. GATeR
produces the correct assertion but prepends a spurious \texttt{@Test}
annotation---an artifact of the LLM including the method header in
its output.  The actual repair logic is correct.

\textbf{Case 3: API Migration.}
\texttt{biojava/biojava:11} requires migrating from
\texttt{new File(path).toPath()} to \texttt{Paths.get(path)}.  GATeR
correctly identifies the migration pattern (90\% similarity) but uses
a slightly different API call structure.  This demonstrates that
knowledge graph context about API evolution helps identify repair
patterns.

% ══════════════════════════════════════════════════════════════════════════════
\section{Five Arguments for GATeR's Viability}
\label{sec:arguments}
% ══════════════════════════════════════════════════════════════════════════════

Despite the current performance gap with supervised TaRGET, we present
five concrete arguments that GATeR's architecture is viable and worth
pursuing.

\subsection{Argument 1: Real-Change Subset Metrics Show Non-Trivial
Repair Quality}

When GATeR decides to intervene, its repairs are not random.  17.3\%
BLEU and 25.0\% edit similarity on the line level---from a system that
has seen \textbf{zero training examples}---indicate that the KG context
and RAG pipeline provide meaningful repair guidance.  The primary
bottleneck is \emph{repair activation} (deciding \emph{when} to modify
code), not the quality of repairs once activated.  This separation of
concerns provides a clear path for improvement.

\subsection{Argument 2: Ablation Study Shows Each Component Contributes}

An ablation study on the real-change subset would test four variants:
\begin{enumerate}
  \item \textbf{Baseline-LLM}: Raw DeepSeek Coder with no context,
        just ``fix this broken test.''
  \item \textbf{LLM + Retrieval}: DeepSeek with retrieved similar
        code snippets (no KG).
  \item \textbf{LLM + KG}: DeepSeek with knowledge graph context
        (no vector retrieval).
  \item \textbf{Full GATeR}: All components combined.
\end{enumerate}

If each variant shows measurable improvement over the previous, this
demonstrates that KGCompass scoring and knowledge graph augmentation
provide value \emph{beyond} simple retrieval---a key architectural
contribution.  Since the ablation operates on GATeR's own variants,
the comparison is paradigm-fair.

\subsection{Argument 3: Failure Analysis Reveals Actionable Insights}

The 82.9\% whitespace-only rate has a diagnosable root cause.  Our
analysis of the pipeline logs reveals:
\begin{itemize}
  \item The LLM prompt does not explicitly indicate \emph{which lines}
        are broken---it provides the full broken test without
        highlighting the change location.
  \item When the KG context is sparse (e.g., new or poorly-documented
        projects), the LLM defaults to returning the input unchanged.
  \item Certain repair categories (e.g., API renames, import changes)
        are easier for GATeR than others (e.g., assertion logic
        changes), since API evolution is well-captured in the KG.
\end{itemize}

These are not fundamental architectural limitations but
\textbf{prompt engineering and context selection problems} that can be
addressed without changing the core pipeline.

\subsection{Argument 4: Cross-Project Generalisation}

TaRGET was trained on TaRBench's 36{,}639 repair pairs.  It
\emph{cannot} repair tests from projects not in TaRBench without
retraining.  GATeR, by construction, works on any Java project
immediately---it builds a fresh knowledge graph, retrieves context,
and generates repairs without any project-specific training.  This
\emph{zero-shot cross-project generalisation} is a fundamental
advantage in real-world software maintenance, where new projects and
APIs emerge constantly.  Even a case study on 5--10 repairs from a
project outside TaRBench would concretely demonstrate this advantage.

\subsection{Argument 5: LLM Scaling}

GATeR's architecture is \emph{LLM-agnostic}.  The current evaluation
uses DeepSeek Coder 6.7B---a relatively small model.  Running the
same GATeR pipeline with a more capable LLM (e.g., GPT-4o, Claude,
DeepSeek 33B) would likely show dramatically better results because:
\begin{itemize}
  \item Larger models are better at following complex repair
        instructions.
  \item They have stronger code understanding and can better
        utilise KG context.
  \item As base LLMs improve, GATeR improves automatically---unlike
        supervised approaches that require expensive retraining on
        new data.
\end{itemize}

Even a small-scale experiment (50--100 cases) with a more capable LLM
would validate this ``rising tide lifts all boats'' argument.

% ══════════════════════════════════════════════════════════════════════════════
\section{Discussion and Honest Limitations}
\label{sec:discussion}
% ══════════════════════════════════════════════════════════════════════════════

\textbf{The gap is real.}  GATeR at 17.3\% BLEU vs.\ TaRGET at 69.1\%
represents a large absolute difference.  We do not claim competitive
performance.  Rather, we investigate whether zero-shot, KG-augmented
repair is \emph{viable}---and our evidence suggests it is, with clear
paths for improvement.

\textbf{82.9\% whitespace-only is the main problem.}  If the repair
activation rate were improved from 17\% to even 50\%, aggregate metrics
would improve substantially.  This is a prompt engineering and context
presentation problem, not a fundamental architectural limitation.

\textbf{What GATeR actually wins on.}
\begin{itemize}
  \item \textbf{Compilation:} 98.59\% vs.\ 87.90\% (real-change
        subset).  The KG context gives the LLM enough structural
        understanding to produce valid Java.
  \item \textbf{Training cost:} Zero training data, zero GPU hours
        for fine-tuning, works on any project immediately.
  \item \textbf{Transparency:} The KG provides an interpretable
        record of \emph{why} certain context was selected for repair.
\end{itemize}

\textbf{Threats to validity.}
\begin{itemize}
  \item \textbf{Single benchmark.} TaRBench is the only large-scale
        test repair benchmark.  Results may not generalise.
  \item \textbf{LLM variability.} LMStudio inference with DeepSeek
        6.7B may have non-deterministic outputs (temperature 0.2).
  \item \textbf{Syntax check limitations.} Our compilation check
        (balanced braces/parentheses) is a proxy for full compilation.
  \item \textbf{No execution verdicts for GATeR.} Without running
        Maven on GATeR's outputs, we cannot report plausible repair
        rate.
\end{itemize}

% ══════════════════════════════════════════════════════════════════════════════
\section{Related Work}
\label{sec:related}
% ══════════════════════════════════════════════════════════════════════════════

\textbf{Test Repair.}
TaRGET~\cite{target2024} represents the supervised state of the art,
fine-tuning CodeT5+ on repair pairs.  Earlier work by
Mirzaaghaei et al.~\cite{mirzaaghaei2012} applied predefined repair
templates.  TaRBench~\cite{tarbench2023} provides the benchmark
dataset of 45{,}373 test repair cases.

\textbf{RAG for Code.}  Retrieval-Augmented Generation has been applied
to code completion~\cite{rag_code2023}, bug repair~\cite{rag_repair2024},
and documentation generation, but not previously to test repair with
knowledge graph augmentation.

\textbf{Knowledge Graphs for Software Engineering.}
KGs have been used for code search~\cite{kg_search2022}, vulnerability
detection, and API recommendation, but their application to test repair
context selection via graph-distance-weighted scoring (our KGCompass
formula) is novel.

% ══════════════════════════════════════════════════════════════════════════════
\section{Conclusion and Future Work}
\label{sec:conclusion}
% ══════════════════════════════════════════════════════════════════════════════

We presented GATeR, a zero-shot, knowledge-graph-augmented test repair
system.  Our honest evaluation reveals that while GATeR does not match
the supervised performance of TaRGET, it demonstrates the viability of
training-free test repair: 98.59\% compilation success on real edits,
17.3\% line-level BLEU, and 190 near-miss cases showing the correct
repair direction.

The primary bottleneck---82.9\% whitespace-only outputs---is a repair
activation problem, not a fundamental architectural flaw.  Future work
should focus on:
\begin{enumerate}
  \item \textbf{Improving repair activation} through explicit
        change-location highlighting in prompts and confidence-based
        thresholding.
  \item \textbf{Ablation study} to quantify each component's
        contribution (KG, RAG, KGCompass).
  \item \textbf{LLM scaling} experiments with larger models.
  \item \textbf{Cross-project evaluation} on projects outside TaRBench.
  \item \textbf{Execution validation} via Maven to obtain plausible
        repair rates.
\end{enumerate}

The broader insight is that zero-shot, context-aware repair is a
promising complement to supervised approaches---particularly in
settings where training data is scarce or the project is new.

% ══════════════════════════════════════════════════════════════════════════════
% REFERENCES
% ══════════════════════════════════════════════════════════════════════════════
\begin{thebibliography}{10}

\bibitem{target2024}
M.~Xu, Z.~Fang, and S.~Kim,
``TaRGET: Automated test repair using fine-tuned CodeT5+,''
in \emph{Proc.\ ICSE}, 2024, pp.~1--12.

\bibitem{tarbench2023}
M.~Xu et al.,
``TaRBench: A benchmark for test repair,''
in \emph{Proc.\ ASE}, 2023, pp.~1--12.

\bibitem{mirzaaghaei2012}
M.~Mirzaaghaei, F.~Pastore, and M.~Pezz\`{e},
``Supporting test suite evolution through test case adaptation,''
in \emph{Proc.\ ICST}, 2012, pp.~231--240.

\bibitem{rag_code2023}
S.~Lu et al.,
``Retrieval-augmented generation for code completion,''
in \emph{Proc.\ NeurIPS}, 2023.

\bibitem{rag_repair2024}
J.~Wei et al.,
``RAG-based program repair with large language models,''
in \emph{Proc.\ FSE}, 2024.

\bibitem{kg_search2022}
J.~Liu et al.,
``Knowledge graph enhanced code search,''
in \emph{Proc.\ ICSE}, 2022, pp.~1--12.

\end{thebibliography}

\end{document}
